\chapter{Diseño de la aplicación}
En este capítulo se describe la aplicación bajo un prisma práctico; mientras que el capítulo de análisis introduce las directivas que se han seguido en cuanto a expectativas de usuario o funcionalidades del sistema, este capítulo las presenta a nivel de diseño, planteando cómo se pueden trasladar las cláusulas propuestas a un plan real. En este capítulo se presenta la arquitectura lógica y física del sistema, expresando su organización y comportamiento. Asimismo, se presentan los tipos de datos planteados en el sistema, exponiendo sus características. Finalmente, se expondrá la interfaz del sistema, su aspecto y funcionalidades.
\section{Arquitectura lógica}
\section{Arquitectura física}
\section{Diagrama de clases, secuencia}
\section{Materialización del modelo de datos}
Explicar cómo django usa el ORM para materializar los datos, cómo funciona a nivel de procesamiento y presentación de la información; describir cómo se ha transformado el ER a través de django para hacer la propia aplicación final, cómo ha afectado el framework a los datos
\section{Interfaces de usuario}
Aquí se hablará sobre todo de Docker, los nodos de computación etc.